%%%%%%%%%%%%%%%%%%%%%%%%%%%%%%%%%%%%%%%%%%%%%%%%%%%%%%%%%%%%
%%%%%%%%%%%%%%%%%%%%%%%%%%%%%%%%%%%%%%%%%%%%%%%%%%%%%%%%%%%%
%%%%%%%%%%%%%%%%%%%%%%%%%%%%%%%%%%%%%%%%%%%%%%%%%%%%%%%%%%%%
%%%%%%%%%%%%%%%%%%%%%%%%%%%%%%%%%%%%%%%%%%%%%%%%%%%%%%%%%%%%
\chapter{Tipos de dados compostos}
\label{chap:type}

Nem todos os problemas podem ser resolvidos com variáveis simples como as primitivas já vistas. Bancos de dados requerem o uso de registros com campos diferentes para armazenar por exemplo o nome, endereço, estado civil etc de cada cliente. Problemas que envolvem matrizes, vetores, listas e sequências requerem o uso de estruturas indexadas, ou seja, que permitam acessar componentes individuais através de índices inteiros. Neste capítulo veremos como criar, modificar e consultar estruturas com tipos de dados compostos, que podem ser homogêneas, compostas por itens de mesmo tipo, ou heterogêneas, compostas por itens de tipos diferentes.


%%%%%%%%%%%%%%%%%%%%%%%%%%%%%%%%%%%%%%%%%%%%%%%%%%%%%%%%%%%%
%%%%%%%%%%%%%%%%%%%%%%%%%%%%%%%%%%%%%%%%%%%%%%%%%%%%%%%%%%%%
\section{Estruturas compostas heterogêneas}

As estruturas compostas heterogêneas são também chamadas de registros. São úteis na criação de bancos de dados e podem armazenar em conjunto dados numéricos e de texto. Abaixo está a sintaxe da criação de um tipo de dados de registro definido pelo usuário.

\NEWLINE
{\tt
struct \PLACEHOLDER{novo tipo}\\
\{\\
\TAB  \PLACEHOLDER{tipo} \PLACEHOLDER{identificador};\\
\TAB  \PLACEHOLDER{tipo} \PLACEHOLDER{identificador};\\
\TAB  ...\\
\}\\
}
\NEWLINE

Abaixo está um exemplo de uso de \STRUCT. Neste exemplo a estrutura {\tt aluno} é definida e são inicializadas duas instâncias, {\tt a} e {\tt b}, de duas maneiras diferentes. Finalmente, seu conteúdo é impresso.

\begin{lstlisting}
struct aluno
{
  char nome[50];
  char matricula[14];
  short int nascDia;
  short int nascMes;
  short int nascAno;
};

int main(void)
{
  struct aluno a = {
    .nome = "John Smith",
    .matricula = "202600001010",
    .nascDia = 31,
    .nascMes = 1,
    .nascAno = 2006,
  };
  struct aluno b = {"Johnson Smith", "202500002020", 1, 12, 2005};

  printf("Nome:       %s\n", a.nome);
  printf("Matricula:  %s\n", a.matricula);
  printf("Nascimento: %02d/%02d/%d\n", a.nascDia, a.nascMes, a.nascAno);
  printf("\n");
  printf("Nome:       %s\n", b.nome);
  printf("Matricula:  %s\n", b.matricula);
  printf("Nascimento: %02d/%02d/%d\n", b.nascDia, b.nascMes, b.nascAno);
}
\end{lstlisting}

Na prática, geralmente se usa \TYPEDEF\ para simplificar a declaração das variáveis. O \TYPEDEF\ cria um alias para um tipo de dados definido pelo usuário, evitando a necessidade da palavra struct a cada declaração de variável desse tipo. A data também pode ser agregada em um outro tipo para simplificar seu uso caso necessário o uso em outras partes do código. Abaixo está a sintaxe usada.
\NEWLINE
{\tt
typedef struct\\
\{\\
\TAB  \PLACEHOLDER{tipo} \PLACEHOLDER{identificador};\\
\TAB  \PLACEHOLDER{tipo} \PLACEHOLDER{identificador};\\
\TAB  ...\\
\} \PLACEHOLDER{novo tipo};\\
} 
\NEWLINE

Abaixo está um exemplo que mostra como o código anterior pode ser modificado para usar \TYPEDEF, além de encapsular a data em um novo tipo.

\begin{lstlisting}
typedef struct
{
  short int dia;
  short int mes;
  short int ano;
} data_t;

typedef struct
{
  char nome[50];
  char matricula[14];
  data_t nasc;
} aluno_t;

int main(void)
{
  aluno_t a = {
    .nome = "John Smith",
    .matricula = "202600001010",
    .nasc = {31, 1, 2006},
  };
  aluno_t b = {"Johnson Smith", "202500002020", 1, 12, 2005};

  printf("Nome:       %s\n", a.nome);
  printf("Matricula:  %s\n", a.matricula);
  printf("Nascimento: %02d/%02d/%d\n", a.nasc.dia, a.nasc.mes, a.nasc.ano);
  printf("\n");
  printf("Nome:       %s\n", b.nome);
  printf("Matricula:  %s\n", b.matricula);
  printf("Nascimento: %02d/%02d/%d\n", b.nasc.dia, b.nasc.mes, b.nasc.ano);
\end{lstlisting}

Tipos criados com \STRUCT\ podem ser passados como parâmetro, retornados e até usados em atribuições, como no exemplo abaixo. Experimente como exercício copiar, modificar e exibir o conteúdo da instância que recebeu a cópia para observar como ela realmente armazena uma nova cópia e não é simplesmente uma referência à instância original.

\begin{lstlisting}
aluno_t alunoCopy(aluno_t src)
{
  aluno_t dst = src;
  return dst;
}

void alunoPrint(aluno_t a)
{
  printf("Nome:       %s\n", a.nome);
  printf("Matricula:  %s\n", a.matricula);
  printf("Nascimento: %02d/%02d/%d\n", a.nascDia, a.nascMes, a.nascAno);
}
\end{lstlisting}

Quando vários registros de mesmo tipo são necessários, podem ser armazenados como posições individuais em um vetor, como no exemplo abaixo. Observe que a inicialização agora não é mais feita exatamente como antes, sendo feita com um \textit{type cast}.

\begin{lstlisting}
  aluno_t arrAluno[10];

  arrAluno[0] = (aluno_t){
    .nome = "John Smith",
    .matricula = "202600001010",
    .nasc = {31, 1, 2006},
  };
  arrAluno[1] = (aluno_t){"Johnson Smith", "202500002020", 1, 12, 2005};
\end{lstlisting}


%%%%%%%%%%%%%%%%%%%%%%%%%%%%%%%%%%%%%%%%%%%%%%%%%%%%%%%%%%%%
%%%%%%%%%%%%%%%%%%%%%%%%%%%%%%%%%%%%%%%%%%%%%%%%%%%%%%%%%%%%
\section{Estruturas compostas homogêneas}

As estruturas compostas homogêneas são chamadas \textit{arrays} e são usadas para representar vetores e matrizes. Vetores são representados por \textit{arrays} unidimensionais e matrizes por \textit{arrays} bidimensionais. O número de dimensões de um \textit{array} é dado pelo número de índices usados para acessar uma posição. Também é possível se necessário ter mais dimensões. O número de posições possíveis de cada dimensão é seu tamanho.

Um vetor com três coordenadas é então representado por um \textit{array} com uma dimensão de tamanho três. Uma matriz com três linhas e quatro colunas é representada por um \textit{array} bidimensional com dimensões de tamanho três e quatro

\textit{Arrays} são sequências de valores do mesmo tipo referenciados por um mesmo nome. Um elemento específico da sequência é acessado por meio de um índice inteiro não negativo. Na linguagem C, os \textit{arrays} começam com o índice zero. Um \textit{array} com {\tt n} posições terá como índices válidos os inteiros de zero a {\tt n-1}.

Uma declaração de \textit{array} tem o seguinte formato:

\NEWLINE
{\tt
\PLACEHOLDER{Tipo de dado} \PLACEHOLDER{Identificador} [ \PLACEHOLDER{Tamanho} ];
}
\NEWLINE

A expressão elementar abaixo funciona como uma variável comum, podendo ser usada em expressões aritméticas, de comparação, atribuição, ou como parâmetro em chamadas de funções.

\NEWLINE
{\tt
\PLACEHOLDER{Identificador} [ \PLACEHOLDER{Posição} ]
}
\NEWLINE

Uma atribuição de valor a uma certa posição do \textit{array}, por exemplo, tem o seguinte formato:

\NEWLINE
{\tt
\PLACEHOLDER{Identificador} [ \PLACEHOLDER{Posição} ] = \PLACEHOLDER{Valor};
}
\NEWLINE


Uma matriz é um \textit{array} bidimensional. Sua declaração é semelhante à de um vetor unidimensional.

\NEWLINE
{\tt
\PLACEHOLDER{Tipo de dado} \PLACEHOLDER{Identificador} [ \PLACEHOLDER{Linhas} ][ \PLACEHOLDER{Colunas} ];
}
\NEWLINE


Para acessar um item individual de uma matriz, usar o formato abaixo:

\NEWLINE
{\tt
\PLACEHOLDER{Identificador} [ \PLACEHOLDER{Linha} ][ \PLACEHOLDER{Coluna} ]      
}
\NEWLINE


%%%%%%%%%%%%%%%%%%%%%%%%%%%%%%%%%%%%%%%%%%%%%%%%%%%%%%%%%%%%
%%%%%%%%%%%%%%%%%%%%%%%%%%%%%%%%%%%%%%%%%%%%%%%%%%%%%%%%%%%%
%%%%%%%%%%%%%%%%%%%%%%%%%%%%%%%%%%%%%%%%%%%%%%%%%%%%%%%%%%%%
\section{Exercícios}

\begin{enumerate}

  \item Crie uma funcão \MAIN:
    \begin{enumerate}
      \item antes dela defina \verb|TAM_MAX| como sendo 10 (dez).
      \item dentro dela declare um vetor de inteiros chamado {\tt vet}, com tamanho máximo igual \verb|TAM_MAX|.
      \item declare uma variável inteira chamada {\tt tamanho}, para armazenar o tamanho ocupado desse vetor. Inicialize-a com zero.
      \item crie um laço \FOR\ que, a cada passo, solicita ao usuário que insira valores para cada posição do vetor, lê e armazena o dado digitado.
      \item em caso de sucesso do \SCANF, incremente {\tt tamanho} e continue. Em caso de erro, interrompa o laço com \BREAK.
      \item depois do laço, imprima o tamanho ocupado do vetor. Note que o tamanho ocupado pode ser diferente do tamanho máximo \verb|TAM_MAX|.
      \item imprima o conteúdo de cada uma das posições do vetor usando a seguinte mensagem: \verb|vet[%d] = %d\n"| substituindo o primeiro \verb|%d| pela posição, e o segundo pelo conteúdo.
    \end{enumerate}

  \item Crie uma funcão \MAIN\ e dentro dela, declare um vetor de reais chamado {\tt numeros}. Seu tamanho máximo é definido por uma macro chamada \verb|TAM_MAX|, que vale 10 (dez). Em seguida, abra um arquivo chamado {\tt input.txt} cujo conteúdo é um valor real em cada linha. Use \FSCANF\ para ler esse arquivo e armazenar os valores no vetor {\tt numeros}, e o tamanho ocupado do vetor na variável inteira chamada {\tt tamanho}. Seu programa deve funcionar para qualquer tamanho de entrada, inclusive para arquivos com mais de dez valores. Nesse caso, o vetor ficará completamente ocupado e o valor de tamanho será 10.

  \item Crie um programa que aloca um \textit{array} com tamanho 4$\times$10. Use o \textit{array} para armazenar os valores abaixo considerando $x$ variando de 1 a 10.
    \begin{enumerate}
      \item Na primeira coluna o valor de $x$.
      \item Na segunda coluna o valor de $x^2$.
      \item Na terceira coluna o valor de $x^3$.
      \item Na quarta coluna o valor de $x^4$.
      \item Imprima o conteúdo do \textit{array}.
    \end{enumerate}

  \item Crie um programa que testa se um certo $n$ é primo. Armazene os primos encontrados em um vetor e use isso para acelerar o processamento. Use um laço para testar se os primos armazenados no vetor dividem $n$.
    \begin{enumerate}
      \item Teste para $n$ = 101.
      \item Solicite para o usuário que insira algum valor de $n$.
      \item Teste para $n$ variando de 2 a 1000 e imprima o total de números primos encontrados.
      \item Teste para $n$ variando de 3 a 9999 pulando os pares e imprima o total de números primos encontrados.
      %\item Use \TIMER\ para cronometrar o tempo de processamento. Compare com a versão do programa do capítulo anterior.
    \end{enumerate}

  \item Arquivos com extensão PGM armazenam imagens do tipo \textit{portable graymap}. São arquivos muito simples, contendo como primeiro byte a letra P seguida de um dígito: 2 indica que as intensidades dos pixels estão em formato de texto e 5 indica formato binário. Linhas começando com {\tt \#} são comentários e devem ser ignoradas. Depois disso vêm três valores inteiros: largura, altura e intensidade máxima dos pixels, que tipicamente é 255 em imagens com 8 bits por amostra. Finalmente vêm os dados dos pixels, armazenados em formato \textit{raster}, ou seja, primeiro a linha superior, da esquerda para a direita e assim por diante. Obtenha ou crie um arquivo PGM. Em seguida, crie uma função \MAIN e dentro dela
    \begin{enumerate}
      \item leia o arquivo e se certifique de que os primeiros bytes são compatíveis com o formato PGM.
      \item leia a largura e a altura da imagem e salve em duas variáveis.
      \item leia o valor da intensidade máxima dos pixels e se certifique de que é 255.
      \item aloque um array de \CHAR\ com tamanho igual a altura $\times$ largura.
      \item leia os valores dos pixels e armazene no array alocado previamente.
    \end{enumerate}

  \item Arquivos com extensão PPM armazenam imagens do tipo \textit{portable pixmap}. São muito semelhantes aos arquivos PGM com algumas diferenças. Contêm como primeiro byte a letra P seguida de um dígito: 3 indica que as intensidades dos pixels estão em formato de texto e 6 indica formato binário. Linhas começando com {\tt \#} são comentários e devem ser ignoradas. Depois disso vêm três valores inteiros: largura, altura e intensidade máxima dos pixels, que tipicamente é 255 em imagens com 8 bits por amostra. Finalmente vêm os dados dos pixels, armazenados em formato \textit{raster}. As amostras de cada pixel são armazenadas na ordem \textit{red}, \textit{green} e \textit{blue}. Obtenha ou crie um arquivo PPM. Em seguida, crie uma função \MAIN e dentro dela
    \begin{enumerate}
      \item leia o arquivo e se certifique de que os primeiros bytes são compatíveis com o formato PPM.
      \item leia a largura e a altura da imagem e salve em duas variáveis.
      \item leia o valor da intensidade máxima dos pixels e se certifique de que é 255.
      \item aloque um array de \CHAR\ com tamanho igual a altura $\times$ largura $\times$ 3 para armazenar os três canais da imagem.
      \item leia os valores dos pixels e armazene no array alocado previamente.
    \end{enumerate}
    
\end{enumerate}
